\section{Introduction}
In recent years, Large Language Models (LLMs) have evolved from simple chatbots into powerful autonomous agents. These modern AI agents can actively plan and complete complex tasks like writing software, managing files, and navigating websites. By using digital tools, these agents act as the ``hands'' of the AI model, allowing it to automate workflows much like a human professional would \cite{agent_survey}. However, despite their potential, current agent systems are very inefficient for everyday use. Today's popular agents lack proper control over how tasks are executed. Because these agents constantly send their entire, uncompressed chat history back and forth to the cloud for every single step, they consume a massive amount of API tokens. This quickly becomes very expensive to operate, especially for long or repetitive tasks.

This financial burden is a widespread problem across the technology industry. Across various professional sectors (such as software development, digital design, search engine optimization (SEO), data analysis, digital marketing, content creation, etc.), there is a strong desire to adopt AI agents for daily workflows. However, the extreme API token costs remain a major barrier to adoption. This issue is not limited to independent professionals or small businesses; recent reports indicate that even large corporations like Uber are now spending more money on AI token consumption than on the human employees the AI was meant to replace \cite{uber_ai_cost}. Furthermore, relying entirely on the cloud makes these agents slow and fully dependent on a continuous internet connection. This poor architectural design makes it difficult for regular users and businesses to use traditional AI agents affordably. Therefore, there is an urgent need for a new type of agent architecture that is smart enough to handle complex tasks, but designed from the ground up to be highly efficient, cost-effective, and capable of running locally.\\


\section{Motivation}
Currently, most AI agents try to solve all types of problems without a clear, structured plan \cite{general_ai_agent1, general_ai_agent2}. This lack of organization causes major issues. Because these systems remember the entire conversation history instead of just the important details, they waste tokens. They also tend to repeat their work for similar tasks instead of reusing past solutions. Furthermore, they usually lack a local supervisor model to help guide their actions efficiently.\\

From a system design perspective, this method is highly wasteful. Instead of saving long conversations, agents should only save the actual steps needed to solve a problem so they can reuse that knowledge later. Adding a local supervisor to manage the agent's tasks can greatly improve speed and reduce unnecessary cloud processing.\\

This research is driven by a simple idea: AI agents should not try to remember everything; they should only remember what matters. Currently, many users waste AI credits and face high costs when using these assistants for daily tasks. By improving how agents manage their memory and adding a local supervisor, this research aims to provide a practical, cost-effective solution to these problems.\\

\section{Problem Statement}
Existing LLM-based agent systems face the following problems:\\
\begin{itemize}
\item \textbf{Unstructured Memory Usage:} Most agents store the full, uncompressed chat history during prolonged interactions. This approach rapidly consumes the available context window, significantly increases token costs, and introduces irrelevant noise that reduces the overall reasoning efficiency of the model\cite{ai_agent-memory}.
\item \textbf{Lack of Workflow Representation:} Tasks are frequently executed in an ad-hoc manner without a structured planning mechanism. This lack of modular representation makes the execution process opaque, which complicates debugging, hinders error recovery, and makes systemic optimization highly difficult.
\item \textbf{Inefficient Tool Selection:} Tools and external APIs are often selected through basic prompt-based guessing rather than a deep semantic understanding of the immediate requirement\cite{efficient_tool_select}. This blind selection process frequently leads to incorrect tool usage, execution failures, and repeated actions that waste computational resources.
\item \textbf{Limited Knowledge Reuse:} Current architectures fail to effectively archive and retrieve previously successful workflows. Because historical execution patterns are not reused, the system is forced to solve identical problems from scratch, causing severe redundant computation and unnecessary delays.
\item \textbf{Absence of a Local Supervisor:} The system relies entirely on the primary cloud-based LLM for both complex reasoning and basic task routing. Without a lightweight local intelligence supervisor to pre-process commands, the agent often makes unnecessary external API calls just to understand simple instructions, which heavily impacts both latency and financial cost.
\end{itemize}

\section{Research Objectives}
The main objectives of this research are to overcome the limitations of current cloud-native systems (such as Claude Code and Google Antigravity)\cite{anthropic_claude_code_works, google_antigravity_ide} and to introduce better token economics \cite{token_economics, ai_tokenomics}:\\
\begin{itemize}
\item \textbf{Eliminate Massive Context Windows:} Remove the dependency on huge, uncompressed context windows by retaining only the most critical information, which aligns with the principles of marginal token allocation \cite{marginal_token}.
\item \textbf{Avoid Repetitive Code Generation:} Stop generating unique and disposable code loops for every single task \cite{claude_code_dive}. Instead, the system will reuse persistent and pre-built local tools to save computational effort.
\item \textbf{Maintain Structured JSON Roadmaps:} Move away from unstructured natural language planning. The goal is to generate deterministic JSON roadmaps that contain clear, step-by-step instructions for task execution.
\item \textbf{Minimize Cloud Communication for Minor Tasks:} Reduce network bandwidth and runtime costs by handling small navigations, routine interactions, and previously solved tasks locally instead of calling the main LLM \cite{efficient_agents}.
\item \textbf{Optimize Skill and Tool Transmission:} Avoid passing lengthy text-based descriptions of Model Context Protocol (MCP) servers or external tools directly into the main prompt to save valuable token space.
\end{itemize}

\section{Core Research Areas}
To achieve these objectives, this research focuses on five primary areas of innovation:\\
\begin{enumerate}
    \item \textbf{Smart Memory Management:} Utilizing a local vector database (Chroma DB) to efficiently build context and manage memory without losing important details or over-compressing data.
    \item \textbf{Workflow Reuse via Specialized Agents:} Reusing past task workflows whenever possible. For example, a specialized Browser Agent can remember previous navigation steps to instantly repeat web tasks. This is achieved by dividing work among specialized agents (such as Coding, Browser, OS, and Testing agents).
    \item \textbf{Persistent Tool Retrieval:} Storing pre-built and previously generated tools inside the vector database so they can be easily retrieved and reused, avoiding the need to generate new code from scratch.
    \item \textbf{Tiny LLMs for Minor Tasks:} Using a very small, lightweight local AI model to handle basic jobs, such as matching past tasks, picking the right tools, and selecting the correct specialized agent.
    \item \textbf{Dynamic Difficulty Routing:} Implementing a smart routing system that calculates a task's difficulty score. While the Tiny LLM is currently in the testing stage, future implementations will route tasks that fail locally to either a low-cost, mid-cost, or high-cost cloud model based on exactly how difficult the problem is.
\end{enumerate}

\section{Research Questions}
This research aims to answer the following questions:\\
\begin{enumerate}
    \item \textbf{RQ1:} How many API tokens does a local-first routing system (local supervisor $\rightarrow$ tiny model $\rightarrow$ cloud) save compared to traditional cloud-only agents?
    \item \textbf{RQ2:} Can using a local database (Chroma DB) and clear, step-by-step JSON plans prevent the AI from wasting memory and repeating the exact same work?
    \item \textbf{RQ3:} How much time and processing power is saved by reusing permanent, pre-made local tools instead of forcing the AI to write temporary new code for every task?
    \item \textbf{RQ4:} How accurately can a fast, local supervisor understand a user's request, pick the right tools, and assign the task without ever asking a cloud AI for help?
\end{enumerate}
